\documentclass[12pt,a4paper]{article}

\usepackage[T2A]{fontenc}
\usepackage[utf8]{inputenc}
\usepackage[russian]{babel}
\usepackage{amsmath,amssymb,amsthm}
\usepackage{array}
\usepackage{booktabs}
\usepackage{enumitem}
\usepackage{tikz}
\usepackage{hyperref}
\usepackage{geometry}
\geometry{margin=2.2cm}

\usetikzlibrary{arrows.meta,positioning,shapes,fit}

\newtheorem{definition}{Определение}
\newtheorem{remark}{Замечание}
\newtheorem{example}{Пример}

\title{Билет №62 \\ \small Организационная структура проекта. Отслеживание и оперативное управление проектом. Сдача. Определение циклической разработки (Round-Trip Engineering), варианты реализации. (СпБГУ)}
\author{Сериков Владислав Александрович}
\date{16 мая 2026}

\begin{document}

\maketitle
\tableofcontents
\newpage

\section{Проект и его организационный контекст}

\begin{definition}
\textbf{Проект} --- это временное предприятие, направленное на создание уникального продукта, услуги или результата. Временность означает наличие начала и завершения, а уникальность означает, что результат проекта отличается от рутинной операционной деятельности.
\end{definition}

Проект существует не изолированно, а внутри некоторой организации. Поэтому его успех зависит не только от технических решений, но и от:
\begin{itemize}
    \item распределения полномочий;
    \item структуры подчинения;
    \item доступности ресурсов;
    \item механизма принятия решений;
    \item правил контроля и отчетности;
    \item способа взаимодействия с заказчиком и пользователями.
\end{itemize}

\begin{definition}
\textbf{Организационная структура проекта} --- это совокупность ролей, полномочий, связей подчинения, каналов коммуникации и правил взаимодействия участников проекта, обеспечивающая достижение целей проекта.
\end{definition}

Организационная структура отвечает на вопросы:
\begin{itemize}
    \item кто принимает решения;
    \item кто отвечает за результат;
    \item кто управляет ресурсами;
    \item кому подчиняются исполнители;
    \item как эскалируются проблемы;
    \item как контролируется ход работ.
\end{itemize}

\section{Основные типы организационных структур проекта}

В проектном управлении обычно выделяют три базовых типа организационных структур:
\begin{enumerate}
    \item функциональная структура;
    \item проектная структура;
    \item матричная структура.
\end{enumerate}

\subsection{Функциональная структура}

\begin{definition}
\textbf{Функциональная организационная структура} --- это структура, в которой сотрудники сгруппированы по профессиональным функциям: разработка, тестирование, аналитика, эксплуатация, маркетинг, финансы и т.д.
\end{definition}

В такой структуре каждый сотрудник обычно имеет одного непосредственного руководителя --- функционального менеджера.

\[
\text{Генеральный директор}
\rightarrow
\begin{cases}
\text{Отдел разработки},\\
\text{Отдел тестирования},\\
\text{Отдел аналитики},\\
\text{Отдел эксплуатации}.
\end{cases}
\]

\subsubsection*{Схема функциональной структуры}

\begin{center}
\begin{tikzpicture}[
    node distance=1.3cm,
    every node/.style={draw, rounded corners, align=center, minimum width=3cm, minimum height=0.8cm},
    arrow/.style={-{Latex}}
]
\node (ceo) {Руководитель\\организации};
\node (dev) [below left=of ceo] {Отдел\\разработки};
\node (test) [below=of ceo] {Отдел\\тестирования};
\node (analyst) [below right=of ceo] {Отдел\\аналитики};

\node (d1) [below=of dev] {Разработчики};
\node (t1) [below=of test] {Тестировщики};
\node (a1) [below=of analyst] {Аналитики};

\draw[arrow] (ceo) -- (dev);
\draw[arrow] (ceo) -- (test);
\draw[arrow] (ceo) -- (analyst);
\draw[arrow] (dev) -- (d1);
\draw[arrow] (test) -- (t1);
\draw[arrow] (analyst) -- (a1);
\end{tikzpicture}
\end{center}

\subsubsection*{Преимущества}
\begin{itemize}
    \item высокая профессиональная специализация;
    \item устойчивость структуры;
    \item простая вертикаль подчинения;
    \item эффективное накопление компетенций внутри отделов;
    \item удобство для повторяющихся процессов.
\end{itemize}

\subsubsection*{Недостатки}
\begin{itemize}
    \item слабая ориентация на конкретный проект;
    \item межфункциональные согласования могут быть медленными;
    \item менеджер проекта имеет ограниченные полномочия;
    \item сотрудники могут воспринимать проект как второстепенную задачу;
    \item возможны конфликты приоритетов между проектами и функциональными задачами.
\end{itemize}

Функциональная структура хорошо подходит для стабильных организаций, где проекты невелики, а основная деятельность носит операционный характер.

\subsection{Проектная структура}

\begin{definition}
\textbf{Проектная организационная структура} --- это структура, в которой команда формируется вокруг конкретного проекта, а руководитель проекта обладает существенными полномочиями по управлению сроками, бюджетом, содержанием и ресурсами.
\end{definition}

В проектной структуре сотрудники часто полностью выделяются в проектную команду.

\subsubsection*{Схема проектной структуры}

\begin{center}
\begin{tikzpicture}[
    node distance=1.3cm,
    every node/.style={draw, rounded corners, align=center, minimum width=3cm, minimum height=0.8cm},
    arrow/.style={-{Latex}}
]
\node (ceo) {Руководитель\\организации};
\node (pm1) [below left=of ceo] {Руководитель\\проекта A};
\node (pm2) [below=of ceo] {Руководитель\\проекта B};
\node (pm3) [below right=of ceo] {Руководитель\\проекта C};

\node (team1) [below=of pm1] {Команда A};
\node (team2) [below=of pm2] {Команда B};
\node (team3) [below=of pm3] {Команда C};

\draw[arrow] (ceo) -- (pm1);
\draw[arrow] (ceo) -- (pm2);
\draw[arrow] (ceo) -- (pm3);
\draw[arrow] (pm1) -- (team1);
\draw[arrow] (pm2) -- (team2);
\draw[arrow] (pm3) -- (team3);
\end{tikzpicture}
\end{center}

\subsubsection*{Преимущества}
\begin{itemize}
    \item сильная ориентация на результат проекта;
    \item высокая скорость принятия решений;
    \item ясная ответственность руководителя проекта;
    \item удобство управления сроками и бюджетом;
    \item высокая вовлеченность команды.
\end{itemize}

\subsubsection*{Недостатки}
\begin{itemize}
    \item возможное дублирование специалистов в разных проектах;
    \item сложность загрузки сотрудников после завершения проекта;
    \item риск изоляции проектных команд;
    \item менее эффективное развитие функциональных компетенций;
    \item высокая зависимость от квалификации руководителя проекта.
\end{itemize}

Проектная структура эффективна для крупных, сложных и критичных проектов, где требуется концентрация ресурсов и быстрое управление.

\subsection{Матричная структура}

\begin{definition}
\textbf{Матричная организационная структура} --- это структура, сочетающая функциональную и проектную организацию. Сотрудник одновременно связан с функциональным руководителем и руководителем проекта.
\end{definition}

В матричной структуре возникает двойное подчинение:
\[
\text{Сотрудник}
\rightarrow
\begin{cases}
\text{Функциональный руководитель},\\
\text{Руководитель проекта}.
\end{cases}
\]

\subsubsection*{Схема матричной структуры}

\begin{center}
\begin{tikzpicture}[
    node distance=1.2cm,
    every node/.style={draw, rounded corners, align=center, minimum width=2.8cm, minimum height=0.75cm},
    arrow/.style={-{Latex}}
]
\node (org) {Руководство};
\node (f1) [below left=of org] {Разработка};
\node (f2) [below=of org] {Тестирование};
\node (f3) [below right=of org] {Аналитика};

\node (p1) [below=2.2cm of f1] {Проект A};
\node (p2) [below=2.2cm of f2] {Проект B};
\node (p3) [below=2.2cm of f3] {Проект C};

\draw[arrow] (org) -- (f1);
\draw[arrow] (org) -- (f2);
\draw[arrow] (org) -- (f3);

\draw[arrow, dashed] (f1) -- (p1);
\draw[arrow, dashed] (f1) -- (p2);
\draw[arrow, dashed] (f2) -- (p1);
\draw[arrow, dashed] (f2) -- (p3);
\draw[arrow, dashed] (f3) -- (p2);
\draw[arrow, dashed] (f3) -- (p3);
\end{tikzpicture}
\end{center}

\subsubsection*{Виды матричной структуры}
\begin{enumerate}
    \item \textbf{Слабая матрица.} Руководитель проекта выполняет скорее координационную роль, а основные полномочия остаются у функциональных руководителей.
    \item \textbf{Сбалансированная матрица.} Полномочия распределены между руководителем проекта и функциональными руководителями.
    \item \textbf{Сильная матрица.} Руководитель проекта обладает значительными полномочиями, близкими к проектной структуре.
\end{enumerate}

\subsubsection*{Преимущества}
\begin{itemize}
    \item эффективное использование специалистов;
    \item сохранение профессиональных функциональных центров;
    \item возможность вести несколько проектов одновременно;
    \item гибкость распределения ресурсов;
    \item сочетание проектной ориентации и функциональной экспертизы.
\end{itemize}

\subsubsection*{Недостатки}
\begin{itemize}
    \item двойное подчинение;
    \item возможные конфликты приоритетов;
    \item повышенные требования к коммуникациям;
    \item необходимость согласования ресурсов;
    \item риск размывания ответственности.
\end{itemize}

\subsection{Сравнение организационных структур}

\begin{center}
\begin{tabular}{|p{3.2cm}|p{3.8cm}|p{3.8cm}|p{3.8cm}|}
\hline
\textbf{Критерий} & \textbf{Функциональная} & \textbf{Матричная} & \textbf{Проектная} \\
\hline
Полномочия руководителя проекта & Низкие & Средние или высокие & Высокие \\
\hline
Доступ к ресурсам & Через функциональных руководителей & По согласованию & Непосредственный \\
\hline
Ориентация на проект & Слабая & Средняя/высокая & Высокая \\
\hline
Профессиональная специализация & Высокая & Высокая & Может ослабляться \\
\hline
Скорость принятия решений & Низкая или средняя & Средняя & Высокая \\
\hline
Риск конфликтов & Между отделами & Между двумя линиями управления & Между проектами за ресурсы \\
\hline
\end{tabular}
\end{center}

\section{Роли в проекте}

Организационная структура проекта конкретизируется через роли.

\subsection{Основные роли}

\begin{itemize}
    \item \textbf{Заказчик} --- сторона, заинтересованная в получении результата проекта.
    \item \textbf{Спонсор проекта} --- лицо или группа лиц, обеспечивающие проект ресурсами и административной поддержкой.
    \item \textbf{Руководитель проекта} --- лицо, ответственное за достижение целей проекта в рамках ограничений по срокам, бюджету, содержанию и качеству.
    \item \textbf{Проектная команда} --- участники, выполняющие работы проекта.
    \item \textbf{Функциональные руководители} --- руководители подразделений, предоставляющих специалистов и компетенции.
    \item \textbf{Пользователи} --- лица, которые будут использовать результат проекта.
    \item \textbf{Стейкхолдеры} --- все заинтересованные стороны, способные повлиять на проект или испытывающие влияние проекта.
\end{itemize}

\subsection{Матрица ответственности RACI}

Для фиксации ответственности часто применяется матрица RACI.

\begin{definition}
\textbf{RACI} --- это метод распределения ответственности, в котором для каждой работы или результата указываются роли:
\begin{itemize}
    \item \textbf{R} --- Responsible: исполняет работу;
    \item \textbf{A} --- Accountable: несет итоговую ответственность;
    \item \textbf{C} --- Consulted: консультирует;
    \item \textbf{I} --- Informed: информируется.
\end{itemize}
\end{definition}

\subsubsection*{Пример RACI-матрицы}

\begin{center}
\begin{tabular}{|p{4cm}|c|c|c|c|}
\hline
\textbf{Работа} & \textbf{PM} & \textbf{Аналитик} & \textbf{Разработчик} & \textbf{Заказчик} \\
\hline
Сбор требований & A & R & C & C \\
\hline
Проектирование архитектуры & A & C & R & I \\
\hline
Реализация модуля & A & I & R & I \\
\hline
Приемочное тестирование & A & C & C & R \\
\hline
Сдача проекта & R & C & C & A \\
\hline
\end{tabular}
\end{center}

\section{Отслеживание и оперативное управление проектом}

\subsection{Понятие отслеживания проекта}

\begin{definition}
\textbf{Отслеживание проекта} --- это регулярный сбор, измерение и анализ информации о фактическом состоянии проекта с целью сравнения его с планом.
\end{definition}

Отслеживаются:
\begin{itemize}
    \item сроки;
    \item бюджет;
    \item объем выполненных работ;
    \item качество;
    \item риски;
    \item проблемы;
    \item загрузка ресурсов;
    \item изменения требований;
    \item удовлетворенность заказчика.
\end{itemize}

\subsection{Оперативное управление проектом}

\begin{definition}
\textbf{Оперативное управление проектом} --- это совокупность регулярных управленческих действий, направленных на удержание проекта в допустимых границах по срокам, стоимости, содержанию и качеству.
\end{definition}

Оперативное управление включает:
\begin{itemize}
    \item постановку краткосрочных задач;
    \item контроль выполнения;
    \item выявление отклонений;
    \item анализ причин отклонений;
    \item принятие корректирующих мер;
    \item перераспределение ресурсов;
    \item управление изменениями;
    \item коммуникацию со стейкхолдерами.
\end{itemize}

\subsection{Цикл оперативного управления}

Оперативное управление удобно рассматривать как замкнутый цикл.

\[
\text{Планирование}
\rightarrow
\text{Исполнение}
\rightarrow
\text{Измерение}
\rightarrow
\text{Анализ}
\rightarrow
\text{Коррекция}
\rightarrow
\text{Планирование}.
\]

\begin{center}
\begin{tikzpicture}[
    node distance=1.6cm,
    every node/.style={draw, rounded corners, align=center, minimum width=2.8cm, minimum height=0.8cm},
    arrow/.style={-{Latex}, thick}
]
\node (plan) {План};
\node (exec) [right=of plan] {Исполнение};
\node (measure) [right=of exec] {Измерение};
\node (analysis) [below=of measure] {Анализ};
\node (correct) [left=of analysis] {Коррекция};
\node (update) [left=of correct] {Обновление\\плана};

\draw[arrow] (plan) -- (exec);
\draw[arrow] (exec) -- (measure);
\draw[arrow] (measure) -- (analysis);
\draw[arrow] (analysis) -- (correct);
\draw[arrow] (correct) -- (update);
\draw[arrow] (update) -- (plan);
\end{tikzpicture}
\end{center}

\subsection{Базовые величины контроля}

Пусть:
\begin{itemize}
    \item $P$ --- плановое значение показателя;
    \item $F$ --- фактическое значение;
    \item $\Delta$ --- отклонение.
\end{itemize}

Тогда:
\[
\Delta = F - P.
\]

Для относительного отклонения:
\[
\delta = \frac{F - P}{P} \cdot 100\%.
\]

\begin{example}
Плановая длительность задачи --- $10$ дней. Фактическая длительность --- $12$ дней. Тогда:
\[
\Delta = 12 - 10 = 2,
\]
\[
\delta = \frac{12 - 10}{10} \cdot 100\% = 20\%.
\]
Задача превысила плановую длительность на $2$ дня, или на $20\%$.
\end{example}

\subsection{Контроль сроков}

Для контроля сроков используются:
\begin{itemize}
    \item календарный план;
    \item диаграмма Ганта;
    \item сетевой график;
    \item критический путь;
    \item контрольные точки;
    \item итерационные планы;
    \item burn-down и burn-up диаграммы.
\end{itemize}

\subsubsection*{Диаграмма Ганта}

Диаграмма Ганта показывает задачи проекта во времени.

\begin{center}
\begin{tabular}{|p{4cm}|c|c|c|c|c|}
\hline
\textbf{Задача} & \textbf{Нед. 1} & \textbf{Нед. 2} & \textbf{Нед. 3} & \textbf{Нед. 4} & \textbf{Нед. 5} \\
\hline
Анализ требований & X & X & & & \\
\hline
Проектирование & & X & X & & \\
\hline
Разработка & & & X & X & \\
\hline
Тестирование & & & & X & X \\
\hline
Сдача & & & & & X \\
\hline
\end{tabular}
\end{center}

\subsection{Метод критического пути}

\begin{definition}
\textbf{Критический путь} --- это наиболее длинная по длительности последовательность зависимых работ проекта, определяющая минимально возможную длительность всего проекта.
\end{definition}

Если задерживается задача на критическом пути, то при прочих равных условиях задерживается весь проект.

\subsubsection*{Алгоритм нахождения критического пути}

\begin{enumerate}
    \item Построить список работ проекта.
    \item Для каждой работы определить длительность.
    \item Для каждой работы определить предшественников.
    \item Построить ориентированный ациклический граф зависимостей.
    \item Выполнить прямой проход:
    \begin{itemize}
        \item найти раннее начало $ES$;
        \item найти раннее окончание $EF$.
    \end{itemize}
    \item Выполнить обратный проход:
    \begin{itemize}
        \item найти позднее окончание $LF$;
        \item найти позднее начало $LS$.
    \end{itemize}
    \item Вычислить резерв времени:
    \[
    Slack = LS - ES = LF - EF.
    \]
    \item Работы с нулевым резервом образуют критический путь.
\end{enumerate}

\subsubsection*{Минимальный пример}

Пусть есть работы:

\begin{center}
\begin{tabular}{|c|c|c|}
\hline
\textbf{Работа} & \textbf{Длительность} & \textbf{Предшественники} \\
\hline
A & 2 & --- \\
\hline
B & 3 & A \\
\hline
C & 4 & A \\
\hline
D & 2 & B, C \\
\hline
\end{tabular}
\end{center}

Возможные пути:
\[
A \rightarrow B \rightarrow D: 2 + 3 + 2 = 7,
\]
\[
A \rightarrow C \rightarrow D: 2 + 4 + 2 = 8.
\]

Следовательно, критический путь:
\[
A \rightarrow C \rightarrow D,
\]
а минимальная длительность проекта равна $8$ единицам времени.

\subsection{Контроль стоимости}

Для контроля стоимости используются:
\begin{itemize}
    \item плановый бюджет;
    \item фактические затраты;
    \item прогноз до завершения;
    \item анализ освоенного объема.
\end{itemize}

Основные показатели:
\[
CV = EV - AC,
\]
где $CV$ --- отклонение по стоимости, $EV$ --- освоенный объем, $AC$ --- фактическая стоимость.

\[
SV = EV - PV,
\]
где $SV$ --- отклонение по срокам, $PV$ --- плановый объем.

Индексы:
\[
CPI = \frac{EV}{AC},
\]
\[
SPI = \frac{EV}{PV}.
\]

Интерпретация:
\begin{itemize}
    \item если $CPI < 1$, проект перерасходует бюджет;
    \item если $CPI > 1$, проект использует бюджет эффективнее плана;
    \item если $SPI < 1$, проект отстает от графика;
    \item если $SPI > 1$, проект опережает график.
\end{itemize}

\begin{example}
Пусть:
\[
PV = 100,\quad EV = 80,\quad AC = 120.
\]
Тогда:
\[
CV = EV - AC = 80 - 120 = -40,
\]
\[
SV = EV - PV = 80 - 100 = -20,
\]
\[
CPI = \frac{80}{120} = 0{,}67,
\]
\[
SPI = \frac{80}{100} = 0{,}8.
\]
Проект одновременно перерасходует бюджет и отстает от плана.
\end{example}

\subsection{Контроль качества}

Качество проекта оценивается по двум направлениям:
\begin{enumerate}
    \item качество процесса выполнения проекта;
    \item качество создаваемого продукта.
\end{enumerate}

Для программного проекта могут использоваться следующие показатели:
\begin{itemize}
    \item количество дефектов;
    \item плотность дефектов;
    \item процент прохождения тестов;
    \item покрытие кода тестами;
    \item время устранения дефектов;
    \item число повторно открытых дефектов;
    \item соответствие требованиям;
    \item стабильность сборок.
\end{itemize}

Плотность дефектов:
\[
D = \frac{N}{S},
\]
где $N$ --- количество обнаруженных дефектов, $S$ --- размер продукта, например в тысячах строк кода или функциональных точках.

\subsection{Управление рисками в оперативном управлении}

\begin{definition}
\textbf{Риск проекта} --- это неопределенное событие или условие, которое в случае наступления оказывает положительное или отрицательное влияние на цели проекта.
\end{definition}

Риск обычно оценивают через вероятность и ущерб:
\[
Exposure = Probability \cdot Impact.
\]

\subsubsection*{Пример}

Если вероятность задержки поставки равна $0{,}3$, а ущерб оценивается в $500000$ рублей, то:
\[
Exposure = 0{,}3 \cdot 500000 = 150000.
\]

Ожидаемая величина риска равна $150000$ рублей.

\subsection{Алгоритм оперативного контроля проекта}

\subsubsection*{Шаги алгоритма}

\begin{enumerate}
    \item Зафиксировать базовый план проекта:
    \begin{itemize}
        \item сроки;
        \item бюджет;
        \item содержание;
        \item контрольные точки;
        \item критерии качества.
    \end{itemize}
    \item Организовать регулярный сбор фактических данных:
    \begin{itemize}
        \item отчеты исполнителей;
        \item данные системы управления задачами;
        \item результаты тестирования;
        \item финансовые данные;
        \item журнал рисков и проблем.
    \end{itemize}
    \item Сравнить фактическое состояние с планом.
    \item Вычислить отклонения.
    \item Определить существенные отклонения.
    \item Выполнить анализ причин.
    \item Выбрать управленческое действие:
    \begin{itemize}
        \item оставить без изменений;
        \item перераспределить ресурсы;
        \item изменить последовательность работ;
        \item согласовать изменение содержания;
        \item пересмотреть сроки;
        \item усилить контроль качества;
        \item эскалировать проблему.
    \end{itemize}
    \item Зафиксировать решение.
    \item Обновить планы и реестры.
    \item Довести изменения до заинтересованных сторон.
\end{enumerate}

\subsubsection*{Иллюстрация алгоритма}

\begin{center}
\begin{tikzpicture}[
    node distance=1.15cm,
    every node/.style={draw, rounded corners, align=center, minimum width=4.3cm, minimum height=0.75cm},
    arrow/.style={-{Latex}, thick}
]
\node (base) {Базовый план};
\node (fact) [below=of base] {Сбор фактических данных};
\node (compare) [below=of fact] {Сравнение с планом};
\node (dev) [below=of compare] {Оценка отклонений};
\node (cause) [below=of dev] {Анализ причин};
\node (action) [below=of cause] {Корректирующее действие};
\node (update) [below=of action] {Обновление планов};

\draw[arrow] (base) -- (fact);
\draw[arrow] (fact) -- (compare);
\draw[arrow] (compare) -- (dev);
\draw[arrow] (dev) -- (cause);
\draw[arrow] (cause) -- (action);
\draw[arrow] (action) -- (update);
\draw[arrow] (update.east) .. controls +(3,0) and +(3,0) .. (fact.east);
\end{tikzpicture}
\end{center}

\section{Сдача проекта}

\subsection{Понятие сдачи проекта}

\begin{definition}
\textbf{Сдача проекта} --- это процесс формального подтверждения того, что результаты проекта созданы, проверены, соответствуют согласованным требованиям и переданы заказчику или эксплуатирующей организации.
\end{definition}

Сдача проекта является частью завершения проекта, но не тождественна ему. Сдача связана прежде всего с передачей результата, а завершение включает также административное закрытие, финансовые операции, архивирование и анализ опыта.

\subsection{Виды сдачи проекта}

\begin{enumerate}
    \item \textbf{Внутренняя сдача} --- результат принимается внутри проектной организации.
    \item \textbf{Внешняя сдача} --- результат передается заказчику.
    \item \textbf{Промежуточная сдача} --- сдается этап, инкремент или версия.
    \item \textbf{Итоговая сдача} --- сдается весь результат проекта.
    \item \textbf{Приемочная сдача} --- заказчик проверяет соответствие результата критериям приемки.
    \item \textbf{Эксплуатационная сдача} --- результат передается в промышленную эксплуатацию.
\end{enumerate}

\subsection{Критерии приемки}

\begin{definition}
\textbf{Критерии приемки} --- это заранее согласованные условия, при выполнении которых результат проекта считается принятым.
\end{definition}

Для программного проекта критериями приемки могут быть:
\begin{itemize}
    \item реализованы все обязательные требования;
    \item пройдены приемочные тесты;
    \item отсутствуют критические дефекты;
    \item подготовлена пользовательская документация;
    \item выполнены требования безопасности;
    \item обеспечена производительность не ниже заданной;
    \item система развернута в целевой среде;
    \item проведено обучение пользователей;
    \item подписан акт приемки.
\end{itemize}

\subsection{Документы при сдаче проекта}

Типовой комплект документов:
\begin{itemize}
    \item акт сдачи-приемки;
    \item отчет о выполнении проекта;
    \item описание поставленного результата;
    \item пользовательская документация;
    \item эксплуатационная документация;
    \item техническая документация;
    \item протоколы испытаний;
    \item список известных ограничений;
    \item инструкции по сопровождению;
    \item реестр изменений;
    \item реестр открытых вопросов;
    \item отчет по рискам;
    \item финансовый отчет;
    \item отчет об извлеченных уроках.
\end{itemize}

\subsection{Алгоритм сдачи программного проекта}

\subsubsection*{Шаги алгоритма}

\begin{enumerate}
    \item Проверить полноту реализации требований.
    \item Выполнить внутреннее тестирование.
    \item Устранить критические и блокирующие дефекты.
    \item Подготовить поставочный комплект.
    \item Развернуть систему в приемочной или промышленной среде.
    \item Провести демонстрацию результата заказчику.
    \item Выполнить приемочные испытания.
    \item Зафиксировать найденные замечания.
    \item Устранить обязательные замечания.
    \item Получить формальное подтверждение приемки.
    \item Передать документацию и права доступа.
    \item Передать проект в сопровождение или эксплуатацию.
    \item Закрыть договорные и административные обязательства.
    \item Провести итоговый анализ проекта.
\end{enumerate}

\subsubsection*{Иллюстрация процесса сдачи}

\begin{center}
\begin{tikzpicture}[
    node distance=1.1cm,
    every node/.style={draw, rounded corners, align=center, minimum width=4.8cm, minimum height=0.75cm},
    arrow/.style={-{Latex}, thick}
]
\node (ready) {Готовность результата};
\node (test) [below=of ready] {Внутреннее тестирование};
\node (package) [below=of test] {Подготовка поставки};
\node (accept) [below=of package] {Приемочные испытания};
\node (fix) [below=of accept] {Устранение замечаний};
\node (sign) [below=of fix] {Подписание акта};
\node (transfer) [below=of sign] {Передача в эксплуатацию};

\draw[arrow] (ready) -- (test);
\draw[arrow] (test) -- (package);
\draw[arrow] (package) -- (accept);
\draw[arrow] (accept) -- (fix);
\draw[arrow] (fix) -- (sign);
\draw[arrow] (sign) -- (transfer);
\end{tikzpicture}
\end{center}

\subsection{Завершение проекта}

Завершение проекта включает:
\begin{itemize}
    \item подтверждение выполнения работ;
    \item закрытие контрактов;
    \item освобождение ресурсов;
    \item передачу результата;
    \item архивирование документации;
    \item закрытие финансовых обязательств;
    \item анализ успехов и ошибок;
    \item фиксацию извлеченных уроков.
\end{itemize}

\subsection{Типичные проблемы при сдаче}

\begin{itemize}
    \item неформализованные критерии приемки;
    \item расхождение ожиданий заказчика и команды;
    \item неполная документация;
    \item наличие критических дефектов;
    \item недостаточное участие пользователей;
    \item изменения требований в конце проекта;
    \item отсутствие готовности эксплуатационной среды;
    \item неясность ответственности за сопровождение.
\end{itemize}

\section{Циклическая разработка и Round-Trip Engineering}

\subsection{Общее понятие циклической разработки}

\begin{definition}
\textbf{Циклическая разработка} --- это подход к созданию программных систем, при котором разработка выполняется в виде повторяющихся циклов, включающих анализ, проектирование, реализацию, тестирование и уточнение результата.
\end{definition}

В отличие от строго последовательной модели, циклическая разработка допускает возвраты к предыдущим этапам. Это особенно важно в программной инженерии, поскольку требования часто уточняются в процессе разработки.

Обобщенный цикл:
\[
\text{Требования}
\rightarrow
\text{Проектирование}
\rightarrow
\text{Реализация}
\rightarrow
\text{Тестирование}
\rightarrow
\text{Анализ}
\rightarrow
\text{Уточнение требований}.
\]

\begin{center}
\begin{tikzpicture}[
    node distance=1.4cm,
    every node/.style={draw, rounded corners, align=center, minimum width=3.2cm, minimum height=0.8cm},
    arrow/.style={-{Latex}, thick}
]
\node (req) {Требования};
\node (design) [right=of req] {Проектирование};
\node (code) [right=of design] {Реализация};
\node (test) [below=of code] {Тестирование};
\node (review) [left=of test] {Анализ};
\node (change) [left=of review] {Уточнение};

\draw[arrow] (req) -- (design);
\draw[arrow] (design) -- (code);
\draw[arrow] (code) -- (test);
\draw[arrow] (test) -- (review);
\draw[arrow] (review) -- (change);
\draw[arrow] (change) -- (req);
\end{tikzpicture}
\end{center}

\subsection{Round-Trip Engineering}

\begin{definition}
\textbf{Round-Trip Engineering, RTE} --- это подход и набор инструментальных механизмов, обеспечивающих взаимную синхронизацию нескольких связанных артефактов разработки, например модели, исходного кода, конфигурационных файлов и документации.
\end{definition}

Наиболее распространенный случай:
\[
\text{UML-модель}
\leftrightarrow
\text{Исходный код}.
\]

Round-Trip Engineering объединяет два направления:
\begin{enumerate}
    \item \textbf{Forward Engineering} --- генерация кода или других артефактов из модели.
    \item \textbf{Reverse Engineering} --- восстановление модели или документации из существующего кода.
\end{enumerate}

\subsection{Forward Engineering}

\begin{definition}
\textbf{Forward Engineering} --- это преобразование более абстрактного описания системы в более конкретную реализацию.
\end{definition}

Например:
\[
\text{Диаграмма классов UML}
\rightarrow
\text{Заготовки классов на Java}.
\]

\subsubsection*{Минимальный пример}

UML-класс:

\[
\boxed{
\begin{array}{l}
\text{User}\\
\hline
-id: int\\
-name: String\\
\hline
+getName(): String
\end{array}
}
\]

Может быть преобразован в код:

\begin{verbatim}
public class User {
    private int id;
    private String name;

    public String getName() {
        return name;
    }
}
\end{verbatim}

\subsection{Reverse Engineering}

\begin{definition}
\textbf{Reverse Engineering} --- это восстановление более абстрактного представления системы по ее реализации.
\end{definition}

Например:
\[
\text{Исходный код}
\rightarrow
\text{Диаграмма классов UML}.
\]

\subsubsection*{Минимальный пример}

Из кода:

\begin{verbatim}
class Order {
    private Customer customer;
    private List<Item> items;
}
\end{verbatim}

можно восстановить модель:
\[
Order \rightarrow Customer,
\]
\[
Order \rightarrow Item.
\]

То есть класс \texttt{Order} связан с классами \texttt{Customer} и \texttt{Item}.

\subsection{Собственно Round-Trip Engineering}

\begin{definition}
\textbf{Round-Trip Engineering в узком смысле} --- это поддержание согласованности между моделью и кодом при изменениях в любом из этих артефактов.
\end{definition}

Схематично:
\[
\text{Модель}
\xrightarrow{\text{генерация}}
\text{Код}
\xrightarrow{\text{анализ}}
\text{Модель}
\xrightarrow{\text{обновление}}
\text{Код}.
\]

\begin{center}
\begin{tikzpicture}[
    node distance=2.2cm,
    every node/.style={draw, rounded corners, align=center, minimum width=3.2cm, minimum height=1cm},
    arrow/.style={-{Latex}, thick}
]
\node (model) {Модель\\UML};
\node (code) [right=of model] {Исходный\\код};

\draw[arrow, bend left=20] (model) to node[above, draw=none] {Forward} (code);
\draw[arrow, bend left=20] (code) to node[below, draw=none] {Reverse} (model);
\end{tikzpicture}
\end{center}

\section{Варианты реализации Round-Trip Engineering}

\subsection{Вариант 1: модель как главный артефакт}

В этом варианте модель является основным источником истины.

\[
\text{Модель} \rightarrow \text{Код}.
\]

Изменения вносятся преимущественно в модель, затем код генерируется автоматически.

\subsubsection*{Преимущества}
\begin{itemize}
    \item высокая архитектурная целостность;
    \item хорошая документация;
    \item удобство анализа на уровне модели;
    \item возможность генерации повторяющегося кода.
\end{itemize}

\subsubsection*{Недостатки}
\begin{itemize}
    \item сложность поддержки генераторов;
    \item риск потери ручных изменений в коде;
    \item модель может быть слишком громоздкой;
    \item не все аспекты реализации удобно моделировать.
\end{itemize}

\subsection{Вариант 2: код как главный артефакт}

В этом варианте основным источником истины является исходный код.

\[
\text{Код} \rightarrow \text{Модель}.
\]

Модель строится автоматически из кода и используется для анализа, документирования и коммуникации.

\subsubsection*{Преимущества}
\begin{itemize}
    \item код всегда актуален;
    \item меньше риск расхождения реализации и документации;
    \item удобно для сопровождения существующих систем;
    \item хорошо подходит для анализа legacy-систем.
\end{itemize}

\subsubsection*{Недостатки}
\begin{itemize}
    \item восстановленная модель может быть слишком детальной;
    \item из кода трудно восстановить исходные проектные намерения;
    \item архитектурные решения могут быть скрыты в деталях реализации.
\end{itemize}

\subsection{Вариант 3: двусторонняя синхронизация}

В этом варианте изменения могут вноситься и в модель, и в код, а инструмент пытается синхронизировать артефакты.

\[
\text{Модель}
\leftrightarrow
\text{Код}.
\]

\subsubsection*{Преимущества}
\begin{itemize}
    \item гибкость;
    \item удобство для команд, где часть участников работает с моделями, а часть --- с кодом;
    \item возможность поддерживать документацию и реализацию согласованными.
\end{itemize}

\subsubsection*{Недостатки}
\begin{itemize}
    \item сложность разрешения конфликтов;
    \item не все изменения имеют однозначное отображение;
    \item требуется дисциплина работы с инструментом;
    \item возможна потеря семантики при преобразовании.
\end{itemize}

\subsection{Вариант 4: частичная генерация}

На практике часто применяется компромиссный подход:
\begin{itemize}
    \item из модели генерируются интерфейсы, каркасы классов, DTO, конфигурации;
    \item бизнес-логика пишется вручную;
    \item ручные участки защищаются от перезаписи.
\end{itemize}

Для этого применяются:
\begin{itemize}
    \item специальные секции пользовательского кода;
    \item шаблоны генерации;
    \item наследование;
    \item partial classes;
    \item аннотации;
    \item соглашения об именовании.
\end{itemize}

\subsection{Вариант 5: синхронизация через метамодель}

В модельно-ориентированной разработке артефакты могут быть представлены как экземпляры некоторой метамодели.

\[
\text{Артефакт}
\rightarrow
\text{Модель}
\rightarrow
\text{Метамодель}.
\]

В таком подходе синхронизация выполняется не напрямую между кодом и диаграммой, а через общее формальное представление.

\section{Алгоритм Round-Trip Engineering}

\subsection{Общий алгоритм двусторонней синхронизации}

\begin{enumerate}
    \item Загрузить текущее состояние модели.
    \item Загрузить текущее состояние исходного кода.
    \item Построить внутреннее представление модели.
    \item Построить внутреннее представление кода, например абстрактное синтаксическое дерево.
    \item Сопоставить элементы модели и кода:
    \begin{itemize}
        \item классы;
        \item интерфейсы;
        \item методы;
        \item поля;
        \item связи;
        \item пакеты;
        \item зависимости.
    \end{itemize}
    \item Найти различия между артефактами.
    \item Классифицировать различия:
    \begin{itemize}
        \item добавление;
        \item удаление;
        \item переименование;
        \item изменение сигнатуры;
        \item изменение связи;
        \item конфликт.
    \end{itemize}
    \item Применить правила синхронизации.
    \item При наличии конфликтов запросить решение пользователя.
    \item Обновить модель.
    \item Обновить код.
    \item Проверить корректность результата.
    \item Сохранить согласованное состояние.
\end{enumerate}

\subsection{Пример синхронизации}

Исходная модель:
\[
User:
\begin{cases}
id: int,\\
name: String.
\end{cases}
\]

Исходный код:

\begin{verbatim}
public class User {
    private int id;
    private String name;
}
\end{verbatim}

Разработчик добавил в код поле:

\begin{verbatim}
private String email;
\end{verbatim}

Инструмент обратного анализа обнаруживает новое поле и обновляет модель:

\[
User:
\begin{cases}
id: int,\\
name: String,\\
email: String.
\end{cases}
\]

Если затем в модели добавляется операция:
\[
+getEmail(): String,
\]
то при прямой генерации в код может быть добавлен метод:

\begin{verbatim}
public String getEmail() {
    return email;
}
\end{verbatim}

\section{Проблемы Round-Trip Engineering}

\subsection{Проблема неоднозначности отображения}

Не всякая конструкция кода имеет очевидный аналог в UML-модели. Например:
\begin{itemize}
    \item лямбда-выражения;
    \item рефлексия;
    \item метапрограммирование;
    \item динамическая генерация объектов;
    \item сложные шаблоны проектирования;
    \item условная компиляция.
\end{itemize}

\subsection{Проблема потери информации}

При преобразовании:
\[
\text{Модель} \rightarrow \text{Код}
\]
может появиться информация, которой не было в модели.

При преобразовании:
\[
\text{Код} \rightarrow \text{Модель}
\]
может потеряться проектный смысл, который не выражен явно в коде.

\subsection{Проблема конфликтов}

Конфликт возникает, если один и тот же элемент был изменен и в модели, и в коде.

\begin{example}
В модели метод был переименован:
\[
calculatePrice() \rightarrow computePrice().
\]

Одновременно в коде разработчик изменил сигнатуру:
\begin{verbatim}
calculatePrice() -> calculatePrice(Currency currency)
\end{verbatim}

Инструмент должен определить, является ли это:
\begin{itemize}
    \item двумя независимыми изменениями;
    \item конфликтом;
    \item переименованием с изменением параметров.
\end{itemize}
\end{example}

\subsection{Проблема уровня абстракции}

Модель обычно должна быть абстрактнее кода. Если в модель автоматически выгрузить все детали реализации, она может стать слишком сложной и потерять ценность как средство понимания.

\section{Связь организационного управления и циклической разработки}

Циклическая разработка влияет на организацию проекта:
\begin{itemize}
    \item планирование становится итеративным;
    \item сдача может выполняться по инкрементам;
    \item контроль строится на коротких циклах обратной связи;
    \item требования уточняются в процессе проекта;
    \item возрастает роль коммуникации с заказчиком;
    \item управление изменениями становится постоянной деятельностью.
\end{itemize}

В традиционном проекте часто используется схема:
\[
\text{План} \rightarrow \text{Исполнение} \rightarrow \text{Сдача}.
\]

В итеративном проекте используется схема:
\[
\left(
\text{Планирование}
\rightarrow
\text{Разработка}
\rightarrow
\text{Проверка}
\rightarrow
\text{Сдача инкремента}
\right)^n.
\]

\section{Минимальный пример организации программного проекта}

Пусть разрабатывается веб-сервис для регистрации пользователей.

\subsection*{Организационная структура}

Выбрана слабая матричная структура:
\begin{itemize}
    \item аналитик выделен из отдела аналитики;
    \item разработчик --- из отдела разработки;
    \item тестировщик --- из отдела качества;
    \item руководитель проекта координирует работы;
    \item функциональные руководители управляют загрузкой сотрудников.
\end{itemize}

\subsection*{Контроль проекта}

Контрольные точки:
\begin{enumerate}
    \item согласованы требования;
    \item реализован API регистрации;
    \item реализована проверка email;
    \item пройдены функциональные тесты;
    \item выполнена приемка заказчиком.
\end{enumerate}

Метрики:
\begin{itemize}
    \item процент выполненных задач;
    \item количество открытых дефектов;
    \item отклонение по срокам;
    \item успешность приемочных тестов.
\end{itemize}

\subsection*{Round-Trip Engineering}

Модель содержит классы:
\[
User,\quad RegistrationService,\quad EmailValidator.
\]

Из модели генерируются заготовки классов. После доработки кода инструмент обновляет диаграмму классов, добавляя новые методы и зависимости.

\section{Краткие выводы}

\begin{enumerate}
    \item Организационная структура проекта определяет распределение полномочий, ресурсов и ответственности.
    \item Функциональная структура удобна для стабильной специализации, но слабо ориентирована на проект.
    \item Проектная структура обеспечивает сильную ориентацию на результат, но может приводить к дублированию ресурсов.
    \item Матричная структура сочетает преимущества функционального и проектного подходов, но требует развитой коммуникации.
    \item Оперативное управление проектом основано на регулярном сравнении факта с планом и принятии корректирующих действий.
    \item Сдача проекта требует заранее определенных критериев приемки и формального подтверждения результата.
    \item Round-Trip Engineering обеспечивает синхронизацию модели и кода.
    \item Основные формы Round-Trip Engineering: генерация кода из модели, восстановление модели из кода и двусторонняя синхронизация.
    \item Главные проблемы Round-Trip Engineering связаны с неоднозначностью преобразований, конфликтами и различием уровней абстракции.
\end{enumerate}

\end{document}
